\section{Resultados}
\label{sec:resultados}

\subsection{Validaci\'on experimental del DA-AFE}

El resultado cuantitativo m\'as s\'olido obtenido hasta la fecha corresponde a la
validaci\'on de la auto-calibraci\'on del \emph{front-end} anal\'ogico, cuyo dise\~no
experimental se describe en la Secci\'on~\ref{sec:met-daafe}.

La Tabla~\ref{tab:results} presenta la media y la desviaci\'on est\'andar muestral del
\emph{offset} residual normalizado sobre diez reinicios apareados. Cada columna es una
configuraci\'on placa--ganancia y cada fila una etapa anal\'ogica en orden de calibraci\'on;
en cada celda, el valor superior corresponde a la referencia fija compartida y el
inferior a la condici\'on calibrada. Leer una columna de arriba hacia abajo muestra c\'omo
el \emph{offset} se reduce etapa por etapa.

\begin{table}[h]
\centering
\scriptsize
\setlength{\tabcolsep}{3.0pt}
\caption{\emph{Offset} de continua residual (\% del fondo de escala del rango nominal de
\SI{5}{\volt} del ADC), media $\pm$ desviaci\'on est\'andar muestral sobre $N=10$ reinicios
apareados. En cada celda: arriba, referencia fija compartida; abajo, condici\'on
calibrada. Menor es mejor.}
\label{tab:results}
\begin{tabular}{lcccccccc}
\toprule
Etapa & A,\,$\times1$ & A,\,$\times4$ & A,\,$\times16$ & A,\,$\times32$ & A,\,$\times50$ & B,\,$\times16$ & B,\,$\times32$ & B,\,$\times50$\\
\midrule
\multirow{2}{*}{PGA}
 & $3{,}5\pm0{,}2$ & $2{,}4\pm0{,}1$ & $2{,}4\pm0{,}1$ & $1{,}0\pm0{,}2$ & $1{,}2\pm0{,}3$ & $2{,}1\pm0{,}1$ & $1{,}0\pm0{,}2$ & $2{,}5\pm0{,}1$\\
 & $3{,}4\pm0{,}3$ & $2{,}1\pm0{,}3$ & $1{,}9\pm0{,}2$ & $0{,}5\pm0{,}2$ & $1{,}0\pm0{,}5$ & $1{,}7\pm0{,}1$ & $0{,}5\pm0{,}2$ & $1{,}8\pm0{,}3$\\
\addlinespace[2pt]
\multirow{2}{*}{BP}
 & $1{,}5\pm2{,}2$ & $4{,}9\pm0{,}3$ & $3{,}8\pm0{,}9$ & $7{,}0\pm0{,}4$ & $1{,}0\pm0{,}7$ & $0{,}66\pm0{,}02$ & $0{,}66\pm0{,}01$ & $0{,}65\pm0{,}01$\\
 & $1{,}2\pm0{,}2$ & $1{,}0\pm0{,}2$ & $1{,}5\pm1{,}5$ & $1{,}0\pm0{,}5$ & $0{,}9\pm0{,}5$ & $0{,}66\pm0{,}02$ & $0{,}66\pm0{,}01$ & $0{,}65\pm0{,}02$\\
\addlinespace[2pt]
\multirow{2}{*}{SUM}
 & $\mathbf{21{,}8\pm1{,}1}^{\dagger}$ & $6{,}7\pm1{,}0$ & $4{,}9\pm1{,}5$ & $7{,}4\pm0{,}3$ & $0{,}4\pm0{,}2$ & $4{,}5\pm0{,}8$ & $8{,}0\pm0{,}8$ & $0{,}9\pm0{,}6$\\
 & $0{,}8\pm0{,}2$ & $1{,}0\pm0{,}3$ & $1{,}4\pm1{,}5$ & $0{,}6\pm0{,}2$ & $\mathbf{0{,}8\pm0{,}6}^{\ddagger}$ & $0{,}6\pm0{,}5$ & $1{,}2\pm0{,}3$ & $0{,}9\pm0{,}6$\\
\addlinespace[2pt]
\multirow{2}{*}{LP}
 & $\mathbf{45{,}8\pm3{,}1}^{\dagger}$ & $\mathbf{42{,}6\pm4{,}2}^{\dagger}$ & $\mathbf{29{,}6\pm6{,}8}^{\dagger}$ & $\mathbf{12{,}4\pm1{,}3}^{\dagger}$ & $4{,}1\pm2{,}1$ & $3{,}0\pm0{,}1$ & $2{,}9\pm0{,}1$ & $2{,}9\pm0{,}1$\\
 & $3{,}2\pm0{,}9$ & $3{,}9\pm0{,}8$ & $2{,}3\pm1{,}4$ & $1{,}7\pm0{,}7$ & $1{,}2\pm0{,}4$ & $0{,}7\pm0{,}1$ & $0{,}4\pm0{,}2$ & $0{,}7\pm0{,}2$\\
\bottomrule
\end{tabular}

\vspace{2pt}
\footnotesize
$^{\dagger}$Referencia fija $\geq\SI{10}{\percent}$ del fondo de escala.\quad
$^{\ddagger}$\'Unica media calibrada superior a su l\'inea base.
\end{table}

\subsubsection{S\'intesis de los resultados}

\begin{itemize}[itemsep=0.2em]
    \item En las ocho configuraciones evaluadas, el \emph{offset} residual a la salida LP
    ---la que alimenta directamente al ADC--- se redujo de
    \SIrange{2.9}{45.8}{\percent} a \SIrange{0.4}{3.9}{\percent} del fondo de escala.
    Las observaciones individuales calibradas se ubicaron entre
    \SIrange{0.1}{5.2}{\percent}.
    \item Una \textbf{prueba de signos de una cola}~\cite{Conover1999} sobre las
    mediciones apareadas de la etapa LP rechaz\'o la hip\'otesis nula de ausencia de
    reducci\'on en \textbf{las ocho configuraciones} (10 de 10 reinicios apareados, con
    $p\approx10^{-3}$ exacto), resultado que sobrevive a la correcci\'on de Bonferroni
    por comparaciones m\'ultiples.
    \item El resultado se obtuvo con una \textbf{\'unica configuraci\'on de controlador}
    para todas las placas y ganancias, sin reajuste espec\'ifico por placa ni por
    ganancia, lo que respalda la robustez del m\'etodo.
    \item Las etapas cuyo residuo de l\'inea base ya se ubica dentro de su banda muerta
    quedan sin modificar por dise\~no, comportamiento visible en la fila BP de la Placa~B.
\end{itemize}

\subsubsection{Temporizaci\'on}

La correcci\'on visible en las trayectorias de calibraci\'on abarca aproximadamente
\SI{7}{\second}, y las observaciones informales de banco indicaron completitud en torno
a los \SI{10}{\second}. Las ventanas de tiempo l\'imite configuradas totalizan
\SI{57.6}{\second} (aproximadamente \SI{60.5}{\second} incluyendo establecimiento y
refinamiento), excluyendo la sobrecarga de conmutaci\'on. Cuando existe una calibraci\'on
almacenada v\'alida s\'olo se ejecutan los intervalos de establecimiento por etapa
(2560 muestras, alrededor de \SI{1.0}{\second} a \SI{2604}{\hertz}, tasa vigente en esa campa\~na).

No se registraron vencimientos de tiempo l\'imite durante los experimentos reportados. El
\'unico observado en banco correspondi\'o a una placa alimentada mientras su ge\'ofono estaba
en movimiento, cuyo transitorio de alterna viola la hip\'otesis cuasi-est\'atica requerida
por la calibraci\'on \emph{foreground}.

\subsubsection{Limitaciones declaradas}

Por transparencia se explicitan las limitaciones del experimento, tal como fueron
reportadas en la comunicaci\'on cient\'ifica:

\begin{itemize}[itemsep=0.15em]
    \item El estudio se limit\'o a \textbf{dos placas ensambladas independientemente} y a
    un orden de calibraci\'on fijo; no establece desempe\~no estad\'istico poblacional.
    \item No se caracteriz\'o la \textbf{deriva de largo plazo} ni el apareamiento entre
    canales a lo largo de un arreglo completo.
    \item La evidencia concierne al \textbf{punto de operaci\'on} de la cadena de
    acondicionamiento, no al desempe\~no s\'ismico de extremo a extremo.
    \item La mejora relativamente modesta en la etapa PGA refleja el criterio de parada
    seleccionado (banda muerta amplia por tener $|K_{\mathrm{PGA}}|=1$) y no
    inestabilidad del controlador. Un diagn\'ostico separado con banda muerta m\'as estrecha
    redujo el residuo LP de \SI{2.98}{\percent} a \SI{0.84}{\percent}, indicando que hay
    margen de mejora por optimizaci\'on de par\'ametros.
\end{itemize}

No obstante, un \emph{offset} de referencia fija como el \SI{21.8}{\percent} medido a la
salida de la etapa sumadora ya consume m\'as de la quinta parte del rango de entrada del
amplificador y del ADC antes de aplicar se\~nal s\'ismica alguna, lo que dimensiona la
relevancia pr\'actica de la correcci\'on.

\subsection{Validaci\'on funcional del sistema de adquisici\'on}

\begin{itemize}[itemsep=0.2em]
    \item \textbf{Auditor\'ia pre-campo}: la matriz de pruebas de aceptaci\'on E1--E18 se
    complet\'o con resultado satisfactorio en todos sus \'items, tras la correcci\'on de los
    defectos cr\'iticos descritos en la Secci\'on~\ref{sec:auditoria}.
    \item \textbf{Captura continua de 10 minutos}: 52\,080 lotes escritos en SD sobre
    52\,080 esperados, 1\,562\,400 muestras transferidas por radio sobre 1\,562\,400
    esperadas, 4\,687\,200 bytes, una \'unica conexi\'on, cero desconexiones y suma
    SHA-256 recomputada coincidente.
    \item \textbf{Regresi\'on por USB}: 33 de 33 verificaciones satisfactorias, cubriendo
    modos crudo y filtrado, escritura y lectura de SD, confirmaciones de comando y
    restauraci\'on de estado.
    \item \textbf{Repetibilidad de extremo a extremo}: 10 de 10 ciclos completos
    satisfactorios, con integridad verificada en cada uno.
    \item \textbf{Robustez de comunicaci\'on}: corte de WebSocket inyectado durante la
    descarga con recuperaci\'on correcta y transferencia completada; detenci\'on en vivo
    validada tanto durante la captura como durante la descarga.
    \item \textbf{Transiciones de configuraci\'on}: 49 de 49 transiciones de rango de ADC
    y factor de decimaci\'on satisfactorias; 30 de 30 ciclos de apagado y recuperaci\'on.
    \item \textbf{Software de procesamiento}: \emph{gate} automatizado de 58
    verificaciones satisfactorio como condici\'on de avance del porteo web.
\end{itemize}

\subsection{Adquisici\'on en campo y procesamiento MASW}

Se realizaron campa\~nas de medici\'on en tres sitios (con dos campa\~nas separadas en uno
de ellos), cuyos conjuntos de datos crudos
est\'an catalogados y versionados con deduplicaci\'on por suma criptogr\'afica. El cat\'alogo
registra del orden de \textbf{22\,000 archivos crudos} (aproximadamente
\SI{28}{\gibi\byte}) junto con los conjuntos procesados correspondientes.

Sobre esos registros el \emph{pipeline} produce:

\begin{itemize}[itemsep=0.15em]
    \item Registros multicanal coherentes visualizables en formato \emph{waterfall} y de
    trazas, con correcci\'on de polaridad y alineaci\'on entre tandas.
    \item Im\'agenes de dispersi\'on en el dominio frecuencia--velocidad de fase, con
    combinaci\'on ponderada de m\'ultiples disparos.
    \item Curvas de dispersi\'on extra\'idas por modo mediante regiones-pol\'igono, editables
    antes de la inversi\'on.
    \item Perfiles $V_S(z)$ obtenidos por inversi\'on multimodal.
\end{itemize}

\nota{\textbf{Pendiente importante a conversar.} Los perfiles $V_S$ obtenidos a\'un no
est\'an contrastados contra un ensayo de referencia (SPT o similar), por lo que en este
documento se reportan como resultado del \emph{pipeline} y no como caracterizaci\'on
geot\'ecnica validada del sitio. Hay que decidir si conviene: (a) presentar un perfil
$V_S$ concreto ya en la Primera Presentaci\'on aclarando expl\'icitamente que la validaci\'on
externa est\'a pendiente, o (b) reportar \'unicamente la capacidad del \emph{pipeline} y
dejar los perfiles para la Defensa Final. Mi sugerencia es (a) con la aclaraci\'on, porque
demuestra que la cadena completa cierra.}

\nota{Faltan por incorporar a este documento las \textbf{figuras}: cadena anal\'ogica,
diagrama de hardware digital, trayectorias de calibraci\'on, \emph{waterfall} de campo e
imagen de dispersi\'on. Todas existen ya en \texttt{docs/Urucom\_2026\_compact/imagenes/}
y en las salidas del \emph{pipeline}; falta elegir cu\'ales entran seg\'un el recorte que
definamos.}
